\section{Fixing a value and disclosing it}
\label{sec:disclosure}

The public binder makes a choice available immediately. Commitment separates
that event into a binding and a later change of audience. We give an explicit
transition and observation semantics for this separation. Its semantic results
do not assume a complete dependent typing calculus for private networks.

\subsection{Bindings, audiences, and views}

A typed binding carries a fresh identity $c$, an owner $r$, a data type $A$,
and one stored value $a:A$. An audience $S_c\subseteq\Roles$ records who may
directly read the value; $r\in S_c$. The identity and declared interface may
be public without making its stored value public.

For deterministic executions, the primitive successful transitions are
\[
 \Bind_r(c:A:=a;S):
       (\rho,\Omega)\longmapsto
       (\rho[c\mapsto(r,A,a,S)],\Omega\cup\{c\}),
 \tag{20}
\]
\[
 \Pub(c;T):
       (\rho,\Omega)\longmapsto
       (\rho[c.S:=T],\Omega\setminus\{c\}),
 \quad c\in\Omega,\quad S_c\subseteq T.
 \tag{21}
\]
Here $\Omega$ tracks unresolved bindings; the owner authorizes (21). The
update in (21) changes the audience field alone. It reads the already stored
value. A later staged-disclosure calculus could use separate capabilities for
several audience expansions; (20)--(21) give a single-resolution interface.

At each choice occurrence $m$, the specification declares an observation map
$o_m:\Hist_m\to O_m$ from legal histories leading to that occurrence to the
actor's view.
The view includes the authorized cells and the actor's retained local
observations and choices. A local policy is a function of this view. If the
choice type is dependent, require a family $\bar A_m:O_m\to\U$ and use
\[
 A_m(h)=\bar A_m(o_m(h)),\qquad
 f_m:\prod_{o:O_m}\bar A_m(o).
 \tag{22}
\]
Thus the type of the demanded answer is itself meaningful from the information
available at that choice.

One simple discipline permits transmitted types to depend only on the current
public context. Owner-local computations may additionally depend on the
owner's private values. Transmitting a type that depends on an unavailable
value requires a specified abstract package; it is not an implicit use of
that private value in another participant's context.

\begin{example}[A sealed pair of choices]
The script
\[
 \Bind_A(c:A:=a;\{A\});\quad
 \Bind_B(d:B:=b;\{B\});\quad
 \Pub(c;\Roles);\quad\Pub(d;\Roles)
 \tag{23}
\]
fixes both payloads before either is opened. Bob may observe that Alice's
binding event occurred, but his binding policy has no direct input for $a$.
Putting $\Pub(c;\Roles)$ before Bob's binding gives a different interface:
Bob may then choose $b$ as a function of $a$. The distinction is one of
available arguments, independent of the eventual objectives.
\end{example}

\begin{proposition}[Binding preservation]
\label{prop:binding}
In any execution of (20)--(21), a live binding's stored value and type remain
unchanged until its resolution. Successful publication of $c$ exposes that
stored value. Two successful openings of the same retained binding identity,
in executions that preserve its original binding, expose equal values.
\end{proposition}
\begin{proof}
The binding transition allocates a fresh identity, so it overwrites no existing
payload. A publication changes only its audience and resolution status. Every
other transition changes a different identity. Induction on the execution
therefore preserves the stored type and value. Publication reads precisely
that value. The equality statement follows by reading the same immutable
binding in the two executions; it does not identify fresh bindings created
under different identities.
\end{proof}

This is an ideal binding interface. A cryptographic implementation must realize
its binding, disclosure, and observation requirements. The abstract transition
rule is not a cryptographic construction.

\subsection{Uniformity and substitution of views}

Define $h\sim_m h'$ when $o_m(h)=o_m(h')$. Dependent values in the following
equations are compared in the common fiber provided by (22).

\begin{proposition}[Observation uniformity]
\label{prop:uniformity}
The decision $f_m(o_m(h))$ is equal at histories $h\sim_m h'$.
\end{proposition}
\begin{proof}
Both decisions are evaluations of the same dependent function at the same
observation, with the same answer type by (22).
\end{proof}

Uniformity refers to the complete declared view. A secret may influence a
separately published message, which may then influence a receiver. Omitting
the secret's name from a local context does not remove that indirect channel.

\begin{definition}[Observation-respecting interface map]
For histories $h$ of one component and a context map $u$ into another, require
\[
 o_i(h)=o_i(h')\quad\Longrightarrow\quad
 o'_i(u(h))=o'_i(u(h'))
 \tag{24}
\]
for every relevant role and occurrence. In a typed version, answer families
are transported along the same map on observations.
\end{definition}

\begin{lemma}[Composition preserves uniformity]
\label{lem:view-composition}
Composing an observation-uniform policy with an interface map satisfying
(24) gives an observation-uniform policy. Such interface maps are closed
under composition and include identities.
\end{lemma}
\begin{proof}
Histories with equal source observations have equal target observations by
(24). The target policy returns equal decisions at those observations.
For two interface maps, apply (24) first to the first map and then to the
second. Identity maps preserve observation equality directly.
\end{proof}

This elementary law is a necessary semantic component of information-aware
substitution. A syntactic dependent cut rule must ensure its hypotheses,
including the transport of the answer types. A proof of the law does not by
itself supply that rule.

\subsection{Public outcomes and common knowledge}

An ideal public observation is stronger than independent, possibly delayed
delivery of several messages. We can state the required epistemic property
directly, using the standard indistinguishability interpretation of knowledge
\cite{HalpernMoses1990,BaltagMossSolecki1999}.

Fix a set $W$ of possible histories at one public stage. Player $i$ cannot
distinguish $h$ and $h'$ when $o_i(h)=o_i(h')$. A proposition is common
knowledge among a group $S$ at $h$ when it holds at every history reachable
from $h$ by a finite chain of these indistinguishability steps for players
in $S$.

\begin{theorem}[Public readout gives common knowledge]
\label{thm:common-knowledge}
Let $p:W\to P$ be a public readout. Suppose that for every $i\in S$,
\[
 o_i(h)=o_i(h')\quad\Longrightarrow\quad p(h)=p(h').
 \tag{25}
\]
Then at every $h\in W$, the fact $p=p(h)$ is common knowledge among $S$.
\end{theorem}
\begin{proof}
Each permitted indistinguishability step preserves $p$ by (25). Induction on
the length of a chain preserves its initial value. Thus every reachable
history satisfies $p=p(h)$, which is the definition of common knowledge.
\end{proof}

For successful publication in (21), include the binding identity and its
opened value in the common public readout. The theorem then establishes
common knowledge of that factual outcome. Its hypotheses describe the ideal
public stage. A runtime where readers have different delayed snapshots must
justify a different observation model or prove an appropriate refinement.

\subsection{Failure, polarity, and the type of an opening}

Resolution can fail or be withheld. With a coproduct available, a total
interface should then return a result such as
\[
 \Unit\; +\;\sum_{a:A}\Opens(c,a).
 \tag{26}
\]
The left injection reports failure. The right injection supplies a value and
evidence $\Opens(c,a)$ that it is the payload of the successful opening of
$c$. Only the successful branch binds $a:A$ for a dependent continuation.
Resolving a pending
obligation does not imply that it was resolved successfully. In particular,
a continuation after a fallible publication cannot receive an unconditional
public variable $a:A$.

The source language of VegasCore provides an operational instance of this
distinction: it stores immutable private commitments, publishes success or
failure results, and tracks unresolved commitments in its syntax indices
\cite{VegasCore2026}. The successful transition model above isolates the
binding and audience laws; it does not include that implementation's retry,
exit, or scheduling policies.

At the endpoint level, publication has a producing side and receiving sides.
This is the familiar connection between polarity and input/output in session
typing \cite{Griffith2016}. At the global level, binding and publication are
two phases, not an involutive pair of logical duals. Both may be initiated by
the same owner. The useful structure is the combination of a role-owned
choice, its endpoint directions, and an explicit change of audience.
